% =====================================================================
%  chapter04_survival.tex — מדריך הישרדות: שאלות נפוצות ותקלות נפוצות
% =====================================================================
\chapter{מדריך הישרדות}
\chaptersubtitle{שאלות נפוצות ותקלות נפוצות}

הפרק הזה בנוי כמו מוקד שירות: שאלה, תשובה קצרה וצעד הבא. החלק הראשון עוסק בתקלות
יומיומיות, והחלק השני במצבים שבהם צריך לחשוב על פנייה לווטרינר.

\section{תקלות יומיומיות}

\faq{הכלב מסרב להיכנס למעלית. מה עושים?}
לא גוררים. מתרגלים כשאין לחץ: עומדים ליד המעלית ומתגמלים, אחר כך נכנסים ויוצאים בלי לנסוע, ורק
אז נוסעים קומה אחת. בינתיים, אם אפשר, משתמשים במדרגות. אם הסירוב מופיע פתאום אצל כלב שתמיד
נכנס בשמחה, כדאי לחשוב גם על כאב או שינוי בריאותי ולהתייעץ עם הווטרינר.

\faq{הכלב נלחץ ליד תנועה סואנת}
מגדילים מרחק: הולכים בצד הפנימי של המדרכה, בוחרים רחובות שקטים יותר (פרק 2) ומתגמלים כל מבט
רגוע לעבר הכביש. מתחילים מגירוי חלש ורחוק, ומתקרבים בהדרגה לאורך ימים ושבועות. זה עובד טוב יותר
מאשר ״לחשוף עד שיתרגל״.

\faq{הכלב מגיב בעוצמה לכלבים זרים}
יוצרים מרחק ושוברים קו ראייה: עוברים מאחורי רכב חונה, חוצים את הכביש או מסתובבים. לא צועקים
ולא מושכים בחוזקה, כי שני הדברים מעלים את המתח. אם התגובות חוזרות, כדאי לפנות למאלף שעובד
בשיטות מבוססות תגמול \cite{avsab2021}.

\faq{לוח הזמנים התפרק: ישיבה התארכה, רכבת התעכבה}
בשביל זה בדיוק קיימת ״תוכנית ב׳״ מפרק 3: איש קשר שיכול להוציא את הכלב, וחלון זמן גמיש. כשחוזרים
מאוחר, טיול קצר ורגוע עדיף על ניסיון ״לפצות״ בריצה ארוכה בחושך.

\faq{הכלב מוצף בסביבה עמוסה}
יוצאים מהסביבה העמוסה לפינה שקטה, נותנים לכלב לרחרח (ריחרוח מרגיע כלבים רבים) ומסיימים את הטיול
מוקדם. בפעם הבאה מגיעים לאותו מקום בשעה פחות עמוסה ולזמן קצר יותר.

\faq{שכחתי שקית, מים או חטיפים}
שקית: בקשו מבעלי כלבים אחרים, רובם ישמחו לעזור. מים: קצרו את הטיול, בחרו מסלול מוצל ועם
ברזייה. בפעם הבאה תעזור ״ערכת הכפילות״ מפרק 3.

\section{כשזה כבר לא יומיומי}

\faq{מזג האוויר התחמם מאוד באמצע הטיול}
עוברים לצל, מאטים ומציעים מים. מקצרים את הטיול וחוזרים הביתה בדרך המוצלת ביותר. כלבים עם אף
קצר, כלבים כבדים ומבוגרים נמצאים בסיכון גבוה יותר למכת חום \cite{hall2020}. איגוד הווטרינרים
האמריקאי מונה בין סימני האזהרה התנשפות חזקה, ריור מרובה, חולשה או חוסר יציבות, הקאה, לשון כהה
במיוחד, פרכוסים ואובדן הכרה \cite{avma-warm}. אם מופיעים סימנים כאלה, פונים מיד לווטרינר.

\faq{הכלב הרים או בלע משהו מהרחוב}
שאריות פיתה, עצם מבושלת, מסטיק או תרופה שמישהו איבד. ברחוב העירוני יש הכול. חלק ממזונות האדם
מסוכנים לכלבים, למשל שוקולד, ענבים וצימוקים, ומוצרים שמכילים קסיליטול כמו חלק מסוגי המסטיקים
ללא סוכר \cite{aspca-foods}. כשיש חשד להרעלה, הנחיית ארגון \en{ASPCA} היא לפנות מיד לווטרינר או
למוקד הרעלות לבעלי חיים \cite{aspca-apcc}. התרשים הבא מסכם את סדר הפעולות.

\begin{figure}[h]
  \centering
  % =====================================================================
%  eight_block_diagram.tex — decision workflow with EXACTLY EIGHT blocks
%  "מה עושים כשהכלב מרים משהו חשוד מהרצפה?"
%  Blocks = the eight \node commands with style `blk`, `startblk`, `endblk` or `decision`
%  (B1..B8). Edge labels are plain text annotations on arrows, not blocks.
%  Layout: row 1 reads right->left (Hebrew order), row 2 left->right, row 3 bottom.
% =====================================================================
\begin{tikzpicture}[
    font=\sffamily\small,
    node distance=1.0cm,
    base/.style={rtlpar,draw,line width=0.9pt,align=center,text width=3.6cm,minimum height=1.25cm,inner sep=4pt},
    blk/.style={base,rounded corners=3mm,fill=white,draw=navy},
    startblk/.style={base,rounded corners=6mm,fill=sage!30!white,draw=sagedark},
    endblk/.style={base,rounded corners=6mm,fill=gold!30!white,draw=golddark},
    urgent/.style={base,rounded corners=3mm,fill=alert!9!white,draw=alert,line width=1.4pt},
    decision/.style={rtlpar,draw=coraldark,fill=coral!18!white,line width=0.9pt,diamond,aspect=1.9,
                     align=center,text width=2.9cm,inner sep=0pt},
    arr/.style={-{Stealth[length=7pt]},line width=1.0pt,draw=navy},
    lbl/.style={font=\sffamily\footnotesize\bfseries,text=coraldark,fill=paper,inner sep=1.5pt},
    num/.style={circle,fill=navy,text=white,font=\sffamily\scriptsize\bfseries,inner sep=1.2pt,minimum size=13pt}
  ]
  % ---- Row 1 (right -> left) ----
  \node[startblk] (B1) at (12.4,0)  {הכלב הרים או לעס משהו חשוד מהרצפה};
  \node[blk]      (B2) at (7.2,0)   {הרחיקו אותו בעדינות מהמקום, אם זה בטוח};
  \node[blk]      (B3) at (2.0,0)   {נסו לזהות מה זה היה: צלמו או שמרו אריזה ושארית};
  % ---- Row 2 ----
  % "ייתכן" = is it POSSIBLE: an unsure answer is a "yes" (safety-first branching)
  \node[decision] (B4) at (2.0,-3.0) {ייתכן שמשהו נבלע, או שיש סימני מצוקה?};
  \node[urgent]   (B6) at (7.2,-3.0) {פנו מיד לווטרינר או למוקד חירום. אל תגרמו להקאה בעצמכם};
  \node[blk]      (B7) at (12.4,-3.0){פעלו בדיוק לפי הנחיות הווטרינר};
  % ---- Row 3 ----
  \node[blk]      (B5) at (2.0,-6.0) {בוודאות לא נבלע דבר: עקבו אחר הכלב בשעות הקרובות};
  \node[endblk]   (B8) at (12.4,-6.0){תעדו את האירוע ושפרו את המניעה בטיול הבא};

  % ---- step numbers (badges, not blocks) ----
  \foreach \b/\n in {B1/1,B2/2,B3/3,B5/5,B6/6,B7/7,B8/8}{\node[num] at (\b.north east) {\n};}
  % west vertex: the 3->4 arrowhead at the top vertex stays visible (DET-1-07)
  \node[num] at (B4.west) {4};

  % ---- arrows ----
  \draw[arr] (B1) -- (B2);
  \draw[arr] (B2) -- (B3);
  \draw[arr] (B3) -- (B4);
  \draw[arr] (B4) -- node[lbl,above=3pt,pos=0.45] {\RLE{כן}} (B6);
  \draw[arr] (B6) -- (B7);
  \draw[arr] (B4) -- node[lbl,left=2pt] {\RLE{לא, בוודאות}} (B5); % left side: clear of badge 5 and connector (DET-1-07)
  \draw[arr] (B5.east) -- node[lbl,above=1pt] {\RLE{הכול תקין}} (B8.west);
  \draw[arr] ([xshift=-0.4cm]B5.north east) |- ([yshift=-0.45cm]B6.south) -- (B6.south);
  \node[lbl,anchor=north] at ([yshift=-0.5cm,xshift=0.3cm]B6.south) {\RLE{הופיע סימן חריג}};
  \draw[arr] (B7) -- (B8);
\end{tikzpicture}

  \caption{מה עושים כשהכלב מרים משהו חשוד מהרצפה? תרשים זרימה בן שמונה שלבים (נבנה בקוד \en{TikZ}).
  התרשים נועד להתמצאות ואינו מחליף ייעוץ וטרינרי.}
  \label{fig:flow}
\end{figure}

\begin{emergencybox}
\begin{itemize}
  \item \textbf{אל תגרמו להקאה} ואל תיתנו תרופה או ״תרופת סבתא״ בלי הנחיה מפורשת של וטרינר.
  \item \textbf{קשיי נשימה, קריסה, פרכוסים, חולשה פתאומית או הקאות חוזרות}: נוסעים מיד למרפאה
        או לבית חולים וטרינרי, ומתקשרים בדרך.
  \item \textbf{מה לקחת:} את האריזה, שארית החומר או תמונה שלו, ומידע על הזמן המשוער שבו זה קרה.
  \item \textbf{מספרי חירום שמורים מראש:} הווטרינר הקבוע ומוקד החירום הקרוב, כבר עכשיו ולא ברגע האמת.
\end{itemize}
\end{emergencybox}
